\chapter*{\center \Large RECOMENDACIONES}
\addcontentsline{toc}{section}{\bfseries RECOMENDACIONES}
\markboth{RECOMENDACIONES}{RECOMENDACIONES}

\begin{enumerate}
  \item \textbf{Conservar la validación multi-semilla como resultado principal.} La corrida nominal vigente usa 15 semillas por variante con huella de calibración fija, horizonte y endpoint explícitos. Para una conclusión más estrecha, añadir intervalos de confianza y más semillas sin reemplazar la mediana, el rango y los modos de falla.

  \item \textbf{Evaluar frente a una población común de adversarios.} EXP-2 ya completó las 90 celdas previstas. La siguiente ampliación debe comparar MAPPO, IPPO y las líneas base frente a los mismos adversarios reservados para evaluación, además de los compañeros de self-play, para reducir la confusión entre la calidad de la terapia y la dificultad del tumor co-entrenado.

  \item \textbf{Ampliar la verificación del atacante.} El benchmark V2 ya incluye cinco reinicios por terapia, el compañero de self-play y cross-play. Una ampliación útil es medir la estabilidad del peor TTP encontrado al variar el presupuesto de entrenamiento y los reinicios; la mejor respuesta aproximada obtenida no garantiza un óptimo global.

  \item \textbf{Mejorar la calibración biológica.} Separar en el reporte qué parámetros provienen de datos, cuáles de literatura y cuáles de normalización. Cuando sea posible, incorporar curvas dosis--respuesta por línea, incertidumbre bootstrap, datos temporales de crecimiento y mediciones de eliminación del fármaco, en lugar de tratar la brecha de IC50 como identificación completa de la dinámica.

  \item \textbf{Implementar la salud del paciente como un módulo validado, no como una penalización arbitraria.} La salud debe representar variables observables y clínicamente interpretables, por ejemplo toxicidad hematológica acumulada, función orgánica y calidad de vida. Cada componente necesita escala, horizonte temporal, umbral y fuente de calibración; la recompensa debe incluir una restricción o presupuesto de toxicidad separado del TTP. Hasta contar con esos datos, la salud debe permanecer como análisis de sensibilidad y no como conclusión clínica.

  \item \textbf{Ampliar la evaluación externa.} Comparar las políticas con control óptimo, búsqueda evolutiva y otros algoritmos de RL; realizar validación fuera de muestra por línea celular; y probar escenarios con variabilidad de parámetros, observación parcial, retrasos de medición y errores de adherencia. Estas pruebas ayudarán a distinguir una política que aprendió una estrategia general de una que explotó una regularidad particular del simulador.

  \item \textbf{Mantener la trazabilidad de resultados.} Cada tabla o figura debe indicar la fecha, la huella de calibración, el número de semillas, el adversario, el horizonte y si el resultado es histórico, parcial o vigente. Los artefactos históricos no deben mezclarse con la corrida calibrada actual.
\end{enumerate}
